\section{Comparison of FE-MSM and SeqGPLVM}
\label{sec:comparison}
\paragraph{FE-MSM}
The Fixed-Effects MSM (FE-MSM) addresses time-invariant unmeasured confounding by incorporating unit-specific fixed effects directly into the propensity score model.

Unit-Specific Sequential Ignorability: This method relies on the assumption that treatment is randomized conditional on the observed history and a time-constant unit-specific effect $\alpha_i$.
Estimation Procedure: We obtain estimates for the propensity parameters $\beta$ and the unit effects $\alpha_i$ via **Conditional Maximum Likelihood**. This approach is justified under **large-$N$, large-$T$ asymptotics, where the bias from noisy fixed-effect estimation (the incidental parameters problem) fades over time.
Weighted Least Squares: Once the fixed-effect weights are constructed, the MSM parameters $\gamma$ are estimated by solving a weighted estimating equation, which reduces to **Weighted Least Squares (WLS)** for our linear specification. We use a **robust (sandwich) variance estimator** to obtain valid Wald confidence intervals.

\paragraph{SeqGPLVM}

In contrast to the point-estimate approach of FE-MSM, the SeqGPLVM captures unmeasured confounding by inferring a posterior distribution $q(Z)$ over a latent substitute $Z_i$.

Handling Circular Dependency: To avoid the "circular dependency" created by using $Z_i$ as a covariate in a separate propensity model, we use the estimated propensities from the SeqGPLVM **directly** to build the IPTW weights.
Multiple Imputation (MI)-Style Procedure: To propagate the latent uncertainty from the GP-based treatment model, we draw $S$ samples $Z^{(s)} \sim q(Z)$. For each draw, we:
    1.  Compute draw-specific propensities and weights $W^{(s)}$.
    2.  Fit the weighted MSM to obtain a draw-specific estimate $\hat{\gamma}^{(s)}$ and robust covariance $V^{(s)}$.
Rubin’s Rules for Pooling: The final point estimate is the average across draws, $\bar{\gamma} = S^{-1}\sum \hat{\gamma}^{(s)}$. The total variance is calculated using **Rubin’s within/between decomposition**, $T = W + (1 + 1/S)B$, which explicitly accounts for the variability induced by the uncertainty in the inferred latent confounder.



Although the seqgplvm can be used in the settings in which the outcome is observed at every time point, 
since the asymptotic properties and consistency of their IPTW-FE estimator 
relies on having less number of N than T (assumption )
Based on the consistency assumption \ref{ass:C-consistency}, we can define 
truncated potential outcomes 

\begin{comment}

\begin{figure}[ht]
\centering
\resizebox{\textwidth}{!}{
\begin{tikzpicture}[
    font=\small,
    % ADD THIS LINE:
    ampersand replacement=\&, 
    arr/.style={-Latex, thick},
    box/.style={
        draw, rounded corners, align=left,
        inner xsep=5mm, inner ysep=3.5mm,
        text width=0.42\linewidth % Changed from 0.44 to 0.42
    },
    header/.style={
        draw, rounded corners, align=center,
        inner xsep=5mm, inner ysep=2.5mm,
        text width=0.42\linewidth, % Changed from 0.44 to 0.42
        font=\bfseries
    },
    stage/.style={
        font=\bfseries,
        align=right,
        text width=0.20\linewidth, % Increased width
        anchor=east                % Added to fix alignment
    },
    midnote/.style={
        draw, rounded corners, fill=gray!10,
        align=center, inner xsep=2.5mm, inner ysep=2.5mm,
        text width=0.16\linewidth
    }
]

% --- Matrix: All '&' replaced with '\&' ---
\matrix (m) [
    matrix of nodes,
    nodes in empty cells,
    row sep=10mm,
    column sep=12mm,
    nodes={box, anchor=north}
] {
    |[header]| FE--MSM \& |[header]| SeqGPLVM \\
    {\textbf{MSM / estimand.}\\
    Choose $k$ and specify an MSM for $Y_i(d_{T-k:T})$.}
    \&
    {\textbf{MSM / estimand.}\\
    Choose $k$ and specify an MSM for $Y_i(d_{T-k:T})$.}\\
    {\begin{minipage}[t]{\linewidth}
        \begin{enumerate}\setlength\itemsep{1pt}\setlength\parskip{0pt}
        \item No interference
        \item Consistency
        \item Positivity
        \item \textbf{Unit-specific ignorability}
        \item \textbf{Truncated Marginal Structural Model}
        \end{enumerate}
    \end{minipage}}
    \&
    {\begin{minipage}[t]{\linewidth}
        \begin{enumerate}\setlength\itemsep{1pt}\setlength\parskip{0pt}
        \item No interference
        \item Consistency
        \item Positivity
        \item \textbf{Time-invariant unmeasured confounding}
        \item \textbf{Conditional Markov model}
        \item 
        \end{enumerate}
    \end{minipage}}\\
    {\textbf{Identification.}\\
    Given (claimed) ignorability, causal effects can be estimated\\
    by MSM/IPTW using propensities conditional on $(X, A, Z)$.}
    \&
    {\textbf{Identification.}\\
    Given (claimed) ignorability, causal effects can be estimated\\
    by MSM/IPTW using propensities conditional on $(X, A, Z)$. } \\
    {\textbf{Propensity estimation.}\\
    Specify and estimate $\pi_{it}(x_t, d_{t-1}; \beta, \alpha_i)$\\
    (FE MLE). }
    \&
    {\textbf{Stage 1: SeqGPLVM.}\\
    Fit GP+CMM latent factor model to infer $Z_i$ and/or compute\\
    $\hat{\pi}_{it}=\hat{p}(A_{it}\mid A_{i,t-1},X_{it},Z_i)$. } \\
    {\textbf{Construct weights.}\\
    IPTW / stabilized weights from $\hat{\pi}_{it}$\\
    (trim if needed). }
    \&
    {\textbf{Construct weights (this thesis).}\\
    Use $\hat{\pi}_{it}$ from SeqGPLVM directly to build IPTW weights.\\
    \emph{(Original: plug $\hat{Z}_i$ into an outcome model.)} } \\
    {\textbf{Outcome stage.}\\
    Fit MSM by weighted regression / estimating equations. }
    \&
    {\textbf{Outcome stage.}\\
    Fit MSM by weighted regression / estimating equations. } \\
};

% --- Row labels ---
\node[stage] at ($(m-2-1.west)+(-12mm,0)$) {Target};
\node[stage] at ($(m-3-1.west)+(-12mm,0)$) {Identification assumptions};
\node[stage] at ($(m-4-1.west)+(-12mm,0)$) {Estimation assumptions};
\node[stage] at ($(m-5-1.west)+(-12mm,0)$) {Treatment\\ model};
\node[stage] at ($(m-6-1.west)+(-12mm,0)$) {Weights};
\node[stage] at ($(m-7-1.west)+(-12mm,0)$) {Outcome};

% --- Vertical arrows ---
\foreach \r/\rp in {2/3,3/4,4/5,5/6,6/7} {
    \draw[arr] (m-\r-1.south) -- (m-\rp-1.north);
    \draw[arr] (m-\r-2.south) -- (m-\rp-2.north);
}

% --- Central "common output" callout ---
%\node[midnote] (pi) at ($(m-5-1.east)!0.5!(m-5-2.west)$) {Outputs\\$\hat{\pi}_{it}$};
%\draw[arr] (m-5-1.east) -- (pi.west);
%\draw[arr] (m-5-2.west) -- (pi.east);

\end{tikzpicture}
}
\caption{Stage-aligned comparison of FE--MSM and SeqGPLVM. The pipelines differ in how propensities are estimated (fixed effects vs. latent factor + GP), but both feed into IPTW and a weighted MSM fit.}
\label{fig:flow_compare_rows_clean}
\end{figure}
\end{comment}